The purpose of this chapter is to define the architecture of the Early Warning Alarm System space segment. Following the top-level mission objectives and requirements, the spacecraft architecture is decomposed into specialized operational subsystems. To pursue persistent observation relevant to Baltic-region threats and the ambitious latency requirement, the design focuses on a modular architecture; detailed coverage and end-to-end latency performance require verification. This approach couples a standardized satellite bus with highly specialized mission payloads optimized for multi-tier ballistic event tracking.

To support regional protection against missile launches that may originate beyond the Baltic region, the system architecture utilizes a dual-tier space segment constellation:
\begin{itemize}
    \item \textbf{High-Altitude Detection Layer (\abbr{GEO}):} Comprises 6 satellites in Geosynchronous Earth Orbit (35,786 km), spaced uniformly at 60 degree intervals along the equator to provide continuous staring detection across mid-and-low latitudes.
    \item \textbf{Low-Altitude Tracking Layer (\abbr{LEO}):} Comprises a Walker Delta constellation of 72 small satellites operating at an altitude of 1,200 km across 8 high-inclination orbital planes, with 9 satellites per plane. This configuration is intended to support high-latitude and polar observations; quantitative coverage and stereoscopic overlap are subject to the constellation simulation results.
\end{itemize}
The subsystem design presented in this section prioritizes ultra-high-sensitivity detection, structural longevity, and rapid, secure data dissemination.

\subsection{Standard Satellite Bus Architecture}
The \abbr{EWAS} space segment is based on a modular satellite bus utilizing flight-proven Commercial Off-The-Shelf (\abbr{COTS}) components to minimize development risks, control baseline procurement costs, and ensure maximum reliability. This standardized platform provides the essential operational infrastructure via the following subsystems:

\begin{description}
    \item[Structure:] Modular aluminium frame designed for launch load endurance and internal component protection. It incorporates localized structural shielding to mitigate radiation degradation (radiation hardening) encountered during extended operations in \abbr{GEO} and persistent \abbr{LEO} orbital environments.
    \item[Electrical Power System:] Deployable solar arrays and high-capacity Lithium-Ion (Li-Ion) battery storage to guarantee uninterrupted power distribution during orbital eclipse phases. The system is scaled to accommodate standard spacecraft operations alongside the transient peak-power demands of the primary computational assets during active threat tracking.
    \item[Thermal Control System:] Passive thermal regulation utilizing Multi-Layer Insulation (\abbr{MLI}), optical surface reflectors, and dedicated heat-pipe radiators. Special architectural isolation is implemented to handle the cryogenic cooling requirements of the infrared sensors to maintain nominal detector sensitivity.
    \item[Propulsion System:] A low-thrust monopropellant or electric propulsion module. Its primary functions are to perform autonomous station-keeping manoeuvres to maintain precise slot positioning within the \abbr{GEO}/\abbr{LEO} constellation grids, and execute orbital phase management.
    \item[Communications Subsystem:] Incorporates high-gain, highly directional antennas coupled with a dedicated hardware cryptographic module to safeguard data integrity. Crucially, the system features Inter-Satellite Links (\abbr{ISL}) to route data across the orbital shells, providing a potential relay route where direct ground contact is unavailable. Command and observation traffic can use suitable ground-station links or supported \abbr{ISL} relay paths. Neither route is assumed to guarantee the 60-second end-to-end requirement without analysis.
\end{description}

\subsubsection[GEO vs LEO Platform Differentiation]{\abbr{GEO} vs \abbr{LEO} Platform Differentiation}
While both platforms share \abbr{COTS}-derived design patterns, their drastically different orbital environments require separate engineering margins. Table~\ref{tab:bus_differentiation} outlines the specific requirements allocated to the \abbr{GEO} and \abbr{LEO} satellite buses.

\begin{table}[h!]
\centering
\caption[Engineering Requirements for GEO and LEO Satellite Bus Configurations]{Engineering Requirements for \abbr{GEO} and \abbr{LEO} Satellite Bus Configurations}
\label{tab:bus_differentiation}
\begin{tabular}{|l|p{6cm}|p{6cm}|}
\hline
\textbf{Subsystem} & \textbf{\abbr{GEO} Bus Specifications} & \textbf{\abbr{LEO} Bus Specifications} \\ \hline
\textbf{Structure} & Heavy radiation shielding for Geostationary environment. & Lightweight composite/aluminium structure optimized for batch launch. \\ \hline
\textbf{\abbr{EPS}} & Large solar arrays, continuous exposure with rare equinox eclipses. & High-cycle Li-Ion batteries to handle 15 orbital eclipses per day. \\ \hline
\textbf{\abbr{ADCS}} & Earth-pointing nadir stabilization with low angular velocity. & High-agility reaction wheels and magnetorquers for fast target cross-cueing. \\ \hline
\textbf{\abbr{TTC} / \abbr{ISL}} & Persistent high-power ground links and multi-channel \abbr{ISL} routing nodes. & Dynamic, steerable \abbr{ISL} antennas to lock onto moving adjacent \abbr{LEO} nodes. \\ \hline
\end{tabular}
\end{table}

\subsection{Mission-Specific Payloads}
\begin{description}
    \item[Global Multi-Tier Infrared (\abbr{IR}) Sensor Suite:] The primary mission payload is an advanced infrared optical assembly tailored to capture the massive thermal energy of ballistic missile plumes during the boost phase against the cold background of space. The sensors operate in complementary Short-Wave (\abbr{SWIR}) and Medium-Wave (\abbr{MWIR}) bands to minimize false alarms and maximize detection probability.
    \item[\abbr{GEO} Observation Layer:] Focuses on global detection using Wide Field of View (\abbr{WFOV}) staring optics. While spatial resolution is low (approx. 1–2 km per pixel), the extreme radiometric contrast of a rocket motor firing acts as an immediate point-source digital trigger.
    \item[\abbr{LEO} Tracking Layer:] Features Narrow Field of View (\abbr{NFOV}) tracking sensors integrated onto dynamic, dual-axis gimbals. Cued by initial high-altitude alerts, these optics mechanically steer their line of sight to lock onto the target footprint, crossing paths with adjacent \abbr{LEO} nodes to establish multi-angle stereoscopic tracking of the missile post-burnout.
\end{description}
\textbf{Key Design Drivers:} High imagery sampling rates for rapid trajectory tracking and precise signal discrimination to isolate missile plumes from natural solar glints or heavy background atmospheric clutter.

\subsection[On-Board Data Handling and Advanced Processing]{On-Board Data Handling and Advanced Processing}
The On-Board Data Handling (\abbr{OBDH}) system manages all data flows and executes mission-critical software. \abbr{EWAS} utilizes a high-performance computational architecture to process information in real-time while on orbit. To handle data transmission reliably under extreme environmental conditions, the internal data architecture implements the \textbf{SpaceWire} protocol for high-rate payload data and a redundant \textbf{\abbr{CAN} bus} for critical platform telemetry.

Key design drivers include:
\begin{itemize}
    \item \textbf{On-board Event Filtering:} Capability to distinguish between missile plumes and background noise using advanced algorithms to reduce false alarm rates.
    \item \textbf{Latency Optimization:} Significant reduction in warning time by generating detection alerts locally, eliminating the need to downlink massive volumes of raw infrared imagery.
    \item \textbf{Secure Data Management:} Relevant onboard processing records and event logs are held in secure temporary spacecraft storage and downlinked to the Ground Segment Data Archiving System for controlled retention and the technical uses permitted by R-004.
\end{itemize}

\subsection[Attitude Determination and Control System]{Attitude Determination and Control System}
Given the high velocity of \abbr{LEO} nodes and the dynamic motion of \abbr{GEO} platforms relative to sudden threats, the Attitude Determination and Control System (\abbr{ADCS}) must continuously manipulate the spacecraft’s spatial orientation to maintain stable ground tracking. 

Key design drivers include:
\begin{itemize}
    \item \textbf{Precision Pointing:} System-level capability to aim the \abbr{IR} optics at targeted coordinates without image blurring. The target pointing stability is maintained below 2 arcseconds per second to ensure clear pixel capture.
    \item \textbf{Stability Management:} Active elimination of high-frequency micro-vibrations and structural jitter that confuse small plume detection algorithms.
    \item \textbf{Autonomous Target Tracking:} Dynamic adjustments to the spacecraft attitude in real-time to match cued geographic coordinates received from the detection layer.
\end{itemize}

\subsection{Spacecraft Subsystem Interactivity and Command Data Flow}
The baseline operating sequence is Ground Segment-directed rather than GEO-initiated cross-cueing. The GEO detection report reaches the Ground Segment first; ground command-and-control functions then select two suitable LEO satellites. A 60-second end-to-end warning-delivery target is retained, but the feasibility of completing this full sequence has not yet been verified.

\begin{enumerate}
    \item \textbf{GEO Constellation Detection and Reporting:}
    \begin{itemize}
        \item The \abbr{GEO} infrared payload detects a candidate thermal event and performs initial filtering and event classification.
        \item The \abbr{GEO} communications subsystem transmits the event location, timestamp and available detection data to the Ground Segment over an available secure link.
    \end{itemize}
    \item \textbf{Ground-Directed LEO Selection and Cueing:}
    \begin{itemize}
        \item Ground Segment functions identify two suitable \abbr{LEO} satellites based on geometry, visibility, operational state and link availability.
        \item Ground Segment command-and-control functions deliver observation commands via a suitable ground station or a supported \abbr{ISL} relay route, depending on availability. Inter-satellite coordination can supplement the observation when needed.
    \end{itemize}
    \item \textbf{LEO Constellation Tracking:}
    \begin{itemize}
        \item The selected \abbr{LEO} satellites receive the tracking requests, point their sensors and acquire complementary observations.
        \item Onboard processing derives a preliminary trajectory estimate from the available stereoscopic observations.
        \item The satellites encrypt and transmit relevant observations and the preliminary estimate to the Ground Segment through an available supported communication path.
    \end{itemize}
    \item \textbf{Ground Segment Processing and Warning:}
    \begin{itemize}
        \item The Ground Station Network securely receives and forwards satellite data, while central Ground Segment Data Processing performs the principal data fusion.
        \item The system generates a refined trajectory estimate and associated uncertainty-based impact area. It evaluates relevance to the Baltic region.
        \item Only after the two \abbr{LEO} tracking datasets are available and processed does the Early Warning System automatically dispatch a preliminary warning to authorized national authority interfaces; no prior operator confirmation is required.
        \item The Mission Control Center supports subsequent operator review, updates or cancellation, and the Ground Segment Data Archiving System stores the applicable records for the purposes permitted by R-004.
    \end{itemize}
\end{enumerate}

The Ground Segment must not assume that the combined GEO downlink, LEO selection, command uplink/relay, dual-satellite observation, data return, central processing and warning delivery necessarily fits within 60 seconds. This remains a critical engineering verification item. 
